\documentclass[10pt,a4paper]{amsart}
\usepackage[margin=19mm]{geometry}
\usepackage{amsmath,amssymb,mathtools}
\usepackage[scr=rsfs,cal=euler]{mathalfa}
\usepackage{xfrac}
\usepackage[unicode,hidelinks,bookmarks=false]{hyperref}
\newcommand{\ie}{i.e.\ }
\newcommand{\cf}{cf.\ }
\newcommand{\resp}{respectively\ }
\newcommand{\Subsection}{\subsection}
\providecommand{\nobreakdash}{\nobreak\hbox{-}}
\setlength{\emergencystretch}{2em}
\multlinegap0pt
\usepackage{amssymb}
\usepackage{mathtools}
\usepackage{xcolor}
\usepackage[shortlabels]{enumitem}
\newtheorem{lemma}{Lemma}
\newtheorem{theorem}{Theorem}
\usepackage{cleveref}
\crefname{lemma}{Lemma}{Lemmas}
\crefname{theorem}{Theorem}{Theorems}
\title{\textbf{New upper bound for multicolor Ramsey numbers}}
\author{Gang Yang, Yaping Mao}
\date{Source: arXiv:2608.01962v1 (CC BY 4.0)}
\begin{document}
\maketitle

\noindent\emph{Source and licence.} \textbf{New upper bound for multicolor Ramsey numbers}. Version arXiv:2608.01962v1 (CC BY 4.0), distributed by arXiv under the Creative Commons Attribution 4.0 International licence, http://creativecommons.org/licenses/by/4.0/, as recorded in the arXiv licence metadata for this exact version and shown on the abstract page https://arxiv.org/abs/2608.01962v1. Full licence notice: \texttt{licenses/arxiv-2608.01962-CC-BY-4.0.txt}.\par\smallskip

\noindent\textit{Editorial scope.} Counted unit: the article's theorem thm:book (source0.tex lines 1147--1191). The article's complete original proof of it is delivered with it (source0.tex lines 1193--1566, one \texttt{proof} environment of 374 lines). No mathematics is written, repaired or re-derived: the statement and the proof are the article's own lines, in the article's own notation, and no retained line is substituted or re-flowed.  Retained uncounted as the source-local context the proof needs: lines 1018--1046.  Nothing was omitted from any retained span, so the package carries no omission record.  No retained line contains a citation, so the wrapper carries no bibliography and no bibliography entry.  No retained line contains a figure environment, an image-inclusion command or drawing code, so no image is delivered and the package declares no asset dependency.  Editorial formatting only, in addition: an AMS \texttt{amsart} preamble that restates the article's own declarations and macros used by the retained lines, namely the environments \texttt{lemma}, \texttt{theorem} on separate counters, together with the standard packages those lines need.  The retained environments print in this document as Lemma 1, Theorem 1 in order of appearance, whereas the article prints the same items with the numbering of its own full document.  The complete native source of the article as distributed in the arXiv e-print for this exact version is bundled unchanged as \texttt{sources/206658/source0.tex}.
\begin{lemma}\label{lem:key}
Let $r\ge2$ and $d\ge3$ be integers, let $0<\beta\le1$ and $C>0$, and
assume that $\mathcal G_r^{(d)}(\beta,C)$ holds. Let
$X,Y_1,\ldots,Y_r$ be nonempty finite vertex sets in an $r$-edge-colored
complete graph, and put
$
 p_i=p_i(X,Y_i)>0
~(i\in[r]).
$
For every choice of positive real numbers $\alpha_1,\ldots,\alpha_r$, there
exist a vertex $x\in X$, an index $\ell\in[r]$, a real number
$\lambda\ge-1$, a nonempty set $X'\subseteq X$, and nonempty sets
$
 Y_i'\subseteq N_i(x)\cap Y_i
~(i\in[r])
$
such that
\begin{equation}\label{eq:key-X}
 |X'|\ge\beta e^{-C(\lambda+1)^{1/d}}|X|,
\end{equation}
\begin{equation}\label{eq:key-selected}
 p_\ell(X',Y_\ell')\ge p_\ell+\lambda\alpha_\ell,
\end{equation}
and, for every $i\in[r]$,
$
 |Y_i'|=p_i|Y_i|,
 p_i(X',Y_i')\ge p_i-\alpha_i.
$
\end{lemma}

\begin{theorem}\label{thm:book}
Let $r\ge2$ and $d\ge3$ be integers. Let $0<\beta\le1$ and $C>0$, and
assume that $\mathcal G_r^{(d)}(\beta,C)$ holds. Let $t,m$ be positive
integers, let $0<p\le1$, and let $\delta,\lambda_0>0$. Set
\[
 L=\log(1/\delta),
 \qquad
 L_p=\log(1/p),
 \qquad
 \rho=\log(r/\beta),
\]
and suppose that
\[
 0<\delta\le\min\{p/4,1/4\},
 \qquad
 \lambda_0\ge\max\{2,6L\},
 \qquad
 t\ge\lambda_0\delta^{-1/(d-1)}.
\]
Define
\[
 \Pi=\frac{3\delta}{p}+\frac{6LL_p}{\lambda_0},
 \qquad
 \Xi=2\rho+4C\lambda_0^{1/d}
      +\frac{12CL}{\lambda_0^{(d-1)/d}}.
\]
Let $X,Y_1,\ldots,Y_r$ be nonempty vertex sets in an $r$-edge-colored
complete graph. Suppose that
\begin{equation}\label{eq:book-density-hyp}
 |N_i(x)\cap Y_i|\ge p|Y_i|
 \qquad(x\in X,\ i\in[r]),
\end{equation}
\begin{equation}\label{eq:book-page-hyp}
 |Y_i|\ge p^{-t}e^{\Pi t}m
 \qquad(i\in[r]),
\end{equation}
and
\begin{equation}\label{eq:book-reservoir-hyp}
 |X|\ge2rt\,e^{rt\Xi}.
\end{equation}
Then there exist an index $i\in[r]$, a color-$i$ clique $T\subseteq X$
with $|T|=t$, and a set $P\subseteq Y_i\setminus T$ with $|P|=m$ such
that every edge between $T$ and $P$ has color $i$. In particular, $(T,P)$
is a color-$i$ $(t,m)$-book.
\end{theorem}

\begin{proof}
We construct the spines one vertex at a time. At time $s$, let $X(s)$ be the
current reservoir, let $Y_i(s)$ be the current page set for color $i$, and
let $T_i(s)$ be the color-$i$ spine constructed so far. Initially,
$
 X(0)=X,
 Y_i(0)=Y_i,$ and $
 T_i(0)=\emptyset.
$
Whenever the relevant sets are nonempty, write
$p_i(s)=p_i\bigl(X(s),Y_i(s)\bigr)$ and $ p_0=\min_{i\in[r]}p_i(0).
$
By \eqref{eq:book-density-hyp}, $p_0\ge p>0$. At every stage at which the
construction is defined, put
\begin{equation}\label{eq:q-alpha}
 q_i(s)=p_i(s)-p_0+\delta,
 \qquad
 \alpha_i(s)=\frac{q_i(s)}{t}.
\end{equation}

We maintain the following structural invariants for every $i\in[r]$:
\begin{equation}\label{eq:book-invariants}
 \begin{aligned}
 &T_i(s)\subseteq X
   \quad\text{and}\quad T_i(s)\text{ is a color-$i$ clique},\\
 &X(s)\subseteq X\cap
   \bigcap_{h=1}^r\bigcap_{v\in T_h(s)}N_h(v),\qquad Y_i(s)\subseteq Y_i\cap\bigcap_{v\in T_i(s)}N_i(v).
 \end{aligned}
\end{equation}
The invariants hold at time $0$. Because the graph is simple, a vertex does
not belong to its own color-$i$ neighborhood. It follows from
\eqref{eq:book-invariants} that $X(s)$ is disjoint from every current spine
and that $T_i(s)\cap Y_i(s)=\emptyset$ for every $i$. Thus a vertex chosen
from $X(s)$ is always a genuinely new spine vertex.

Suppose that $|T_i(s)|<t$ for every $i\in[r]$ and that the current state is
defined with $p_i(s)>0$ and $q_i(s)>0$ for every $i$. Apply
\Cref{lem:key} to the current reservoir and page sets with parameters
$\alpha_i(s)$. We obtain a vertex $x\in X(s)$, a color $\ell\in[r]$, a
number $\lambda\ge-1$, a set $X'\subseteq X(s)$, and sets
$
 Y_i'\subseteq N_i(x)\cap Y_i(s)
~(i\in[r]).
$
We distinguish two cases.

If $\lambda\le\lambda_0$, then the edges from $x$ to
$X'\setminus\{x\}$ have one of the $r$ colors. Hence some $j\in[r]$
satisfies
\[
 |N_j(x)\cap X'|\ge\frac{|X'|-1}{r}.
\]
We extend the color-$j$ spine by setting
\[
 T_j(s+1)=T_j(s)\cup\{x\},
 \qquad
 X(s+1)=N_j(x)\cap X',
 \qquad
 Y_j(s+1)=Y_j',
\]
and set $T_i(s+1)=T_i(s)$ and $Y_i(s+1)=Y_i(s)$ for $i\ne j$.

If $\lambda>\lambda_0$, we perform a \emph{boost step} in color $\ell$ and
set
$
 X(s+1)=X'$
and $Y_\ell(s+1)=Y_\ell'$,
while leaving every spine and every other page set unchanged.

The structural invariants in \eqref{eq:book-invariants} are preserved. For a
boost step, the spines remain unchanged, $X(s+1)\subseteq X(s)$, and only
$Y_\ell(s)$ is replaced by a subset. For a spine-extension step in color
$j$, the induction hypothesis gives
\[
 x\in X(s)\subseteq\bigcap_{v\in T_j(s)}N_j(v),
\]
so $T_j(s)\cup\{x\}$ is a color-$j$ clique. Moreover,
\[
 X(s+1)\subseteq N_j(x)\cap X(s),
 \qquad
 Y_j(s+1)\subseteq N_j(x)\cap Y_j(s),
\]
and all other spines and page sets remain unchanged. Thus, in either type of
step,
\[
 X(s+1)\subseteq X(s),
 \qquad
 Y_i(s+1)\subseteq Y_i(s)
 \quad(i\in[r]).
\]
We next establish quantitative estimates showing, in particular, that none
of the sets or parameters needed above can vanish before a spine reaches
size $t$.

Formally, carry out the procedure only as long as the current reservoir is
nonempty, and stop at the first state in which some spine has size $t$. If a
proposed spine-extension step were to produce an empty new reservoir, stop
provisionally at that step. The estimates below apply to every legal finite
initial segment. The reservoir estimate will show that the provisional
stopping event is impossible.

\medskip
\noindent\emph{Density bookkeeping.}
For $i\in[r]$, let $B_i(s)$ be the set of earlier boost steps in color $i$,
and let $\lambda(q)$ denote the threshold used at step $q$. Also set
$ e_i(s)=|T_i(s)|$.
Restricting the first argument in a minimum relative degree cannot decrease
that degree. Thus, if the color-$j$ spine is extended at step $s$, then
\[
 p_j(s+1)\ge p_j(X',Y_j')\ge p_j(s)-\alpha_j(s),
\]
whereas, for $i\ne j$, the page set is unchanged and
$X(s+1)\subseteq X(s)$, so $p_i(s+1)\ge p_i(s)$. Similarly, if a boost step
is performed in color $\ell$, then
\[
 p_\ell(s+1)\ge p_\ell(s)+\lambda(s)\alpha_\ell(s),
 \qquad
 p_i(s+1)\ge p_i(s)\quad(i\ne\ell).
\]
Consequently, if the color-$i$ spine is extended at step $s$, then
\[
 q_i(s+1)
 \ge q_i(s)-\alpha_i(s)
 =\left(1-\frac1t\right)q_i(s).
\]
If a boost step is performed in color $i$, then
\[
 q_i(s+1)
 \ge q_i(s)+\lambda(s)\alpha_i(s)
 =\left(1+\frac{\lambda(s)}t\right)q_i(s).
\]
In every other case, $q_i(s+1)\ge q_i(s)$. Iterating these one-step
inequalities gives
\[
 q_i(s)
 \ge q_i(0)\left(1-\frac1t\right)^{e_i(s)}
 \prod_{q\in B_i(s)}\left(1+\frac{\lambda(q)}t\right).
\]
Since $q_i(0)\ge\delta$ and $e_i(s)\le t$ up to and including the first
time a spine reaches size $t$,
\begin{equation}\label{eq:q-product}
 q_i(s)
 \ge\delta\left(1-\frac1t\right)^t
 \prod_{q\in B_i(s)}\left(1+\frac{\lambda(q)}t\right).
\end{equation}
The assumptions imply $t\ge\lambda_0\ge2$, and hence
$(1-1/t)^t\ge1/4$. Therefore
\begin{equation}\label{eq:p-lower}
 q_i(s)\ge\frac\delta4,
 \qquad
 p_i(s)=p_0-\delta+q_i(s)
 \ge p_0-\frac{3\delta}{4}
 \ge p-\frac{3\delta}{4}>0.
\end{equation}
In particular, $q_i(s)>0$, so every $\alpha_i(s)$ in
\eqref{eq:q-alpha} is positive.

Since every relative density is at most $1$ and $\delta\le1/4$, we also have
$q_i(s)\le1+\delta\le5/4$. Combining this with \eqref{eq:q-product} yields
\[
 \prod_{q\in B_i(s)}\left(1+\frac{\lambda(q)}t\right)
 \le\frac5\delta.
\]
Because $\delta\le1/4$, we have $\log(5/\delta)\le3L$, and hence
\begin{equation}\label{eq:boost-log-sum}
 \sum_{q\in B_i(s)}
 \log\left(1+\frac{\lambda(q)}t\right)
 \le3L.
\end{equation}
At every boost step, $\lambda(q)>\lambda_0$. Since $t\ge\lambda_0$ and
$\log(1+x)\ge x/2$ for $0\le x\le1$, each summand in
\eqref{eq:boost-log-sum} is at least
\[
 \log\left(1+\frac{\lambda_0}{t}\right)
 \ge\frac{\lambda_0}{2t}.
\]
Consequently,
\begin{equation}\label{eq:boost-count}
 |B_i(s)|\le\frac{6L}{\lambda_0}t\le t.
\end{equation}

\medskip
\noindent\emph{Page-set retention.}
The page set $Y_i(s)$ is replaced only when the color-$i$ spine is extended
or when a boost step is performed in color $i$. On each such occasion,
\Cref{lem:key} gives
\[
 |Y_i(s+1)|=p_i(s)|Y_i(s)|.
\]
Set
$\xi=\frac{3\delta}{4p}.$
Then $0<\xi\le3/16$, and \eqref{eq:p-lower} gives
$
 p_i(s)\ge p-\frac{3\delta}{4}=p(1-\xi).
$
Let $b_i(s)=|B_i(s)|$ and put $u_i(s)=e_i(s)+b_i(s)$. Up to and including
the first state in which some spine reaches size $t$,
\[
 e_i(s)\le t,
 \qquad
 b_i(s)\le\frac{6L}{\lambda_0}t,
 \qquad
 u_i(s)\le2t.
\]
It follows that
$
 |Y_i(s)|
 \ge p^{t+6Lt/\lambda_0}(1-\xi)^{2t}|Y_i(0)|.
$
Since $\log(1-\xi)\ge-2\xi$ for $0\le\xi\le3/16$, we obtain
\begin{align}
 |Y_i(s)|
 &\ge p^t
 \exp\!\left(-\frac{6LL_p}{\lambda_0}t\right)
 \exp\!\left(-\frac{3\delta}{p}t\right)|Y_i(0)|=p^te^{-\Pi t}|Y_i(0)|.
 \label{eq:page-retention}
\end{align}
By \eqref{eq:book-page-hyp}, every page set occurring in the construction
therefore has size at least $m$.

\medskip
\noindent\emph{The cost of boost steps.}
Suppose that a boost step occurs at time $s$ in color $i$. Then
$
 p_i(s+1)\ge p_i(s)+\lambda(s)\alpha_i(s).
$
Since $p_i(s+1)\le1$ and, by \eqref{eq:p-lower},
$\alpha_i(s)=q_i(s)/t\ge\delta/(4t)$, we have
\begin{equation}\label{eq:lambda-upper}
 \lambda(s)\le\frac{4t}{\delta}.
\end{equation}
For every threshold $\lambda>\lambda_0$ arising from a boost step, we claim
that
\begin{equation}\label{eq:root-over-log}
 \frac{(\lambda+1)^{1/d}}
 {\log(1+\lambda/t)}
 \le\frac{4t}{\lambda_0^{(d-1)/d}}.
\end{equation}
If $\lambda\le t$, then $\lambda\ge\lambda_0\ge2$ and
$\log(1+\lambda/t)\ge\lambda/(2t)$. Hence
\[
 \frac{(\lambda+1)^{1/d}}{\log(1+\lambda/t)}
 \le2t\frac{(\lambda+1)^{1/d}}{\lambda}
 \le4t\lambda^{-(d-1)/d}
 \le\frac{4t}{\lambda_0^{(d-1)/d}}.
\]
If $\lambda>t$, then $\log(1+\lambda/t)\ge\log2$. Moreover,
\eqref{eq:lambda-upper}, $t\ge1$, and $\delta\le1$ give
$\lambda+1\le\frac{5t}{\delta}.
$
The hypothesis $t\ge\lambda_0\delta^{-1/(d-1)}$ is equivalent to
$t^{1/d}\delta^{-1/d}
 \le t\lambda_0^{-(d-1)/d}.
$
Since $d\ge3$ and $5^{1/d}/\log2<4$, inequality
\eqref{eq:root-over-log} follows in this case as well.

Multiplying \eqref{eq:root-over-log} by the corresponding logarithms, using
\eqref{eq:boost-log-sum}, and summing over all colors gives
\begin{equation}\label{eq:boost-root-sum}
 \sum_{q\in B(s)}(\lambda(q)+1)^{1/d}
 \le\frac{12rLt}{\lambda_0^{(d-1)/d}},
 \qquad
 B(s)=\bigcup_{i=1}^rB_i(s).
\end{equation}

\medskip
\noindent\emph{Reservoir retention.}
Set
\[
 \varepsilon_0=\frac\beta r
 \exp\!\left(-C(\lambda_0+1)^{1/d}\right).
\]
Then $0<\varepsilon_0\le1$. In a spine-extension step,
$\lambda\le\lambda_0$, so \Cref{lem:key} and the choice of the majority
color give
\begin{align*}
 |X(s+1)|
 \ge\frac{|X'|-1}{r}\ge\frac\beta r
       e^{-C(\lambda_0+1)^{1/d}}|X(s)|-\frac1r\ge\varepsilon_0|X(s)|-1.
\end{align*}
In a boost step, \Cref{lem:key} gives
$
 |X(s+1)|
 \ge\beta e^{-C(\lambda(s)+1)^{1/d}}|X(s)|.
$
Since $\varepsilon_0\le\beta$, this implies
$
 |X(s+1)|
\ge\varepsilon_0e^{-C(\lambda(s)+1)^{1/d}}|X(s)|.
$
For each step $q$, define
\[
 a_q=
 \begin{cases}
  \varepsilon_0,&\text{if a spine is extended at step $q$},\\
  \varepsilon_0e^{-C(\lambda(q)+1)^{1/d}},
  &\text{if a boost step is performed at step $q$},
 \end{cases}
\]
and let $\epsilon_q=1$ in the first case and $\epsilon_q=0$ in the second.
With $R_q=|X(q)|$, both estimates take the form
$
 R_{q+1}\ge a_qR_q-\epsilon_q,$ where $
 0<a_q\le1.
$
Before or at the first state in which some spine reaches size $t$, fewer than
$rt$ spine extensions and at most $rt$ boost steps have occurred. Indeed, at
such a state the number of spine extensions is at most
$
 t+(r-1)(t-1)=rt-r+1<rt,
$
and \eqref{eq:boost-count} bounds the number of boost steps by $rt$. Thus
there are at most $2rt$ factors $a_q$. Furthermore,
$\log(1/\varepsilon_0)
=\rho+C(\lambda_0+1)^{1/d}
\le\rho+2C\lambda_0^{1/d}.
$
Together with \eqref{eq:boost-root-sum}, this gives, for every constructed
initial segment and also for a putative final spine-extension step that would
empty the reservoir,
\begin{align*}
 \prod_{q=0}^{s-1}a_q
 &=\varepsilon_0^s
 \exp\!\left(-C\sum_{q\in B(s)}(\lambda(q)+1)^{1/d}\right)\ge\exp\!\left(
-2rt\bigl(\rho+2C\lambda_0^{1/d}\bigr)
 -\frac{12CrLt}{\lambda_0^{(d-1)/d}}
 \right)\ge e^{-rt\Xi}.
\end{align*}
Iterating the recurrence for $R_q$ yields
\[
 R_s\ge
 \left(\prod_{q=0}^{s-1}a_q\right)R_0
 -\sum_{j=0}^{s-1}\epsilon_j
  \prod_{q=j+1}^{s-1}a_q.
\]
Every product in the sum is at most $1$, and fewer than $rt$ of the
$\epsilon_j$ are nonzero. Therefore, by
\eqref{eq:book-reservoir-hyp},
$
 |X(s)|=R_s
 \ge e^{-rt\Xi}|X(0)|-rt
 \ge rt>0.
$

We may now close the induction implicit in the construction. The last
reservoir estimate rules out a proposed spine extension with an empty new
reservoir. At every stage before a spine reaches size $t$,
\eqref{eq:p-lower} gives $p_i(s)>0$ and $q_i(s)>0$,
\eqref{eq:page-retention} gives $|Y_i(s)|\ge m\ge1$, and the reservoir
estimate gives $|X(s)|>0$. Therefore \Cref{lem:key} is always applicable,
and the construction cannot stop prematurely.

By \eqref{eq:boost-count}, at most $rt$ boost steps can occur in total. If
every spine had size at most $t-1$, fewer than $rt$ spine extensions could
occur. Since every step is of one of these two types, the process cannot
continue indefinitely without producing a spine of size $t$. Hence, for some
color $i$, the construction reaches a state $s$ with $|T_i(s)|=t$.

Set $T=T_i(s)$. By \eqref{eq:book-invariants}, $T\subseteq X$ is a
color-$i$ clique, and
\[
 Y_i(s)\subseteq Y_i\cap\bigcap_{v\in T}N_i(v).
\]
The page-retention estimate \eqref{eq:page-retention} gives
$|Y_i(s)|\ge m$. Choose a set $P\subseteq Y_i(s)$ with $|P|=m$. Since the
graph is simple, the preceding inclusion implies $Y_i(s)\cap T=\emptyset$.
Consequently,
$
 T\subseteq X$ and
 $P\subseteq Y_i\setminus T,$
and every edge between $T$ and $P$ has color $i$. Thus $(T,P)$ has all the
properties asserted in the theorem.
\end{proof}

\end{document}
